\section{The entire first interior shell}

Write $a$ for the prescribed independence number. Our main extension is the following.

\begin{theorem}[First interior shell, D]\label{thm:shell}
For every integer $k\ge3$ and every integer $2\le a\le k+1$,
\[
f(a+k+1,a,k)=\left\lceil\frac{k(a+k+1)}2\right\rceil.
\]
\end{theorem}

The restriction $a\le k+1$ is exactly $n\ge2a$ on this shell.
Also $n\le ka$, because $a(k-1)\ge k+1$ for $a\ge2$, $k\ge3$.
Every $k$-connected graph has minimum degree at least $k$, so the claimed value is a lower bound by the degree sum.
We give constructions attaining it with exact independence number $a$.

\subsection{Two elementary residual-graph lemmas}

\begin{lemma}[Matching bound]\label{lem:matching}
A nonempty graph of minimum degree at least $b\ge0$ has a matching of size at least $\lceil b/2\rceil$.
\end{lemma}
\begin{proof}
Take a maximal matching $Q$. If a vertex $v$ is unmatched, all its neighbors are matched endpoints, so $2|Q|\ge\dd(v)\ge b$.
If every vertex is matched, $2|Q|$ is the order of the graph and is at least $b+1$.
A subset of $Q$ realizes any smaller prescribed size.
\end{proof}

Number a cycle of order $m$ cyclically by $0,\ldots,m-1$.
Its $r$th power $C_m^r$ joins pairs at cyclic distance at most $r$.
We use $r\ge1$ and $m\ge2r+4$.
The following are classical Harary constructions \cite{Harary}; we include the needed connectivity proof.

\begin{lemma}[Cyclic gaps]\label{lem:harary}
The graph $C_m^r$ is $2r$-connected and $2r$-regular. The following augmentations are $(2r+1)$-connected:
\begin{enumerate}
\item For $m=2q$, add the antipodal matching $\{i,i+q\}$, $0\le i<q$.
The result is $(2r+1)$-regular.
\item For $m=2q+1$, add $\{i,i+q\}$, $0\le i\le q$.
The result has degree $2r+2$ at $q$ and degree $2r+1$ elsewhere.
\end{enumerate}
\end{lemma}
\begin{proof}
In the cyclic list of surviving vertices, consecutive survivors are adjacent if their intervening gap has fewer than $r$ deleted vertices.
With at most one larger gap, these edges form a spanning path or cycle.
Disconnection therefore requires two gaps, each with at least $r$ deletions, proving $2r$-connectivity of the cycle power.
Its degree gives the matching upper bound on connectivity.

For either odd-degree augmentation, only exactly $2r$ deletions require attention.
If the cycle power disconnects, these deletions must form exactly two blocks of length $r$, leaving two nonempty surviving arcs.
Let the shorter arc have length $p$. Rotate coordinates for the geometric argument so that
\[
A=\{0,\ldots,p-1\},\qquad B=\{p+r,\ldots,m-r-1\}.
\]
For $m=2q$, we have $p\le q-r$, and for every $x\in A$,
\[
p+r\le q\le x+q\le q+p-1\le2q-r-1.
\]
Its antipode therefore lies in $B$, yielding a surviving added edge.
For $m=2q+1$, again $p\le q-r$; both $x+q$ and $x+q+1$ lie in $B$, since their indices range from $q$ to $q+p\le2q-r$.
Every vertex in the odd-order construction has an added neighbor at one of these two almost-antipodal positions.
This property survives coordinate rotation even though the chosen edge cover need not be rotation invariant.
Thus an added edge joins the two internally connected arcs in either parity.
The endpoint lists in the odd-order construction overlap only at $q$, giving its degree pattern; the even-order added edges are a perfect matching.
\end{proof}

We will also use
\begin{equation}\label{eq:cycle-alpha}
\alpha(C_m^r)\le\left\lfloor\frac{m}{r+1}\right\rfloor.
\end{equation}
Indeed, the forward cyclic distances between consecutive members of an independent set of size at least two are all at least $r+1$, and sum to $m$.
A singleton also satisfies the bound.

\subsection{An injective missing-neighbor lemma}

\begin{lemma}[Residual lift]\label{lem:lift}
Let $k\ge a\ge3$, $d=k-a$, and let $H$ have order $k+1$.
Choose $M\subseteq V(H)$ with $|M|=a$. Suppose that $\alpha(H)\le a$, that
\[
\dd_H(t)\ge
\begin{cases}d+1,&t\in M,\\d,&t\notin M,\end{cases}
\]
and, when $d\ge1$, that $H$ is $d$-connected.
Add an independent set $S$ of size $a$, put it in bijection with $M$, and join each $s\in S$ to every vertex of $H$ except its assigned vertex.
Then the resulting graph $G$ has $\alpha(G)=a$ and $\kappa(G)=k$.
\end{lemma}
\begin{proof}
Each vertex of $S$ has degree $k$.
With $r_t=1$ for $t\in M$ and $r_t=0$ otherwise,
\[
\dd_G(t)=\dd_H(t)+a-r_t\ge k.
\]
An independent set within either part has size at most $a$.
A mixed independent set contains at most one residual vertex, since a vertex of $S$ misses only one vertex of $H$.
It contains at most one vertex of $S$, since the misses are distinct.
Its size is therefore at most two. The set $S$ witnesses equality $\alpha(G)=a$.

Delete $X$ with $|X|<k$, and write $A=S\setminus X$, $B=V(H)\setminus X$.
Always $|B|\ge2$.
If $A$ is empty, $d=0$ is impossible and fewer than $d=k-a$ residual vertices were deleted, so $d$-connectivity of $H$ suffices.

Suppose $A\ne\varnothing$ and $|B|\ge3$.
Any two vertices of $A$ share a neighbor in $B$, so all of $A$ and every member of $B$ adjacent to it lie in a common component $C$.
At most one residual vertex $t$ can lie outside $C$; it is missed by all of $A$, so $|A|\le r_t$.
If it is disconnected, all its $H$-neighbors have been deleted. Hence
\[
|X|\ge a-|A|+\dd_H(t)\ge a-r_t+\dd_H(t)=\dd_G(t)\ge k,
\]
a contradiction.

If $|B|=2$, precisely $k-1$ residual vertices were deleted, and all of $S$ survives.
Since its $a\ge3$ vertices miss distinct residual vertices, some $s\in S$ is adjacent to both members of $B$.
Every other member of $S$ is adjacent to at least one of them, so the remainder is connected.
Thus $\kappa(G)\ge k$, while the degree-$k$ vertices in $S$ give equality.
\end{proof}

\subsection{The generic residual construction}

Set $m=k+1=a+d+1$, $n=a+m$, $\varepsilon=kn\bmod2$, and $r=\lfloor d/2\rfloor\ge1$.
We construct a $d$-connected residual graph $H$ with degree $d+1$ at exactly $a+\varepsilon$ distinct vertices and degree $d$ elsewhere.
Every case contains $C_m^r$.

If $d=2r$, start with $C_m^r$.
Its complement is $a$-regular, and $\varepsilon$ is the parity of $a$ because $n=2a+d+1$ is odd.
Lemma~\ref{lem:matching} supplies a complement matching of size $(a+\varepsilon)/2$.
Adding it raises exactly $a+\varepsilon$ degrees by one.

If $d=2r+1$ and $m$ is even, start with the antipodal augmentation in Lemma~\ref{lem:harary}.
It is $d$-regular, its complement is $a$-regular, $a$ is even, and $\varepsilon=0$.
Add a complement matching of size $a/2$.

If $d=2r+1$ and $m=2q+1$ is odd, start with the odd-order augmentation.
Its unique higher-degree vertex is $w=q$.
Here $a$ is odd and $\varepsilon=0$.
Every other vertex has degree $a$ in the complement; after deleting $w$, that complement has minimum degree at least $a-1$.
Add a matching of size $(a-1)/2$ in it.
Its endpoints avoid $w$, so precisely $a$ vertices now have degree $d+1$.
All three constructions preserve the starting graph's $d$-connectivity.

By \eqref{eq:cycle-alpha} and $d\le2r+1$,
\[
m=a+d+1\le a+2r+2<(a+1)(r+1),\qquad \alpha(H)\le a.
\]
The strict inequality follows from $r(a-1)-1\ge1$.
Choose any $a$ higher-degree vertices for $M$ and apply Lemma~\ref{lem:lift}.
Every vertex in $S$ and every selected residual vertex has degree $k$.
The residual vertices of degree $d$ also have total degree $k$.
When $\varepsilon=1$, the sole unselected higher-degree vertex has total degree $k+1$.
Consequently
\[
\e(G)=\frac{kn+\varepsilon}{2}=\left\lceil\frac{kn}{2}\right\rceil,
\]
proving Theorem~\ref{thm:shell} in this range.

\subsection{The remaining parameters}

For $a=k$, take a residual matching of $\lceil k/2\rceil$ edges on $k+1$ vertices and choose $k$ distinct endpoints for $M$.
Its independence number is $k+1-\lceil k/2\rceil\le k$.
Lemma~\ref{lem:lift} applies with $d=0$, without requiring residual connectivity.
For even $k$, all degrees in $G$ are $k$; for odd $k$, the single unselected matching endpoint has total degree $k+1$ and every other degree is $k$.
Thus the required edge bound is attained.

For $a=k-1$ and $k\ge4$, use $H=P_{k+1}$ and let $M$ be its $k-1$ internal vertices.
The path is connected, its selected degrees are two, its remaining degrees are one, and
\[
\alpha(P_{k+1})=\left\lceil\frac{k+1}{2}\right\rceil\le k-1.
\]
Lemma~\ref{lem:lift} with $d=1$ gives a $k$-regular extremal graph.
The stronger complete tree characterization from version 0.01 is retained in Theorem~\ref{thm:main} below.
Its $k=3$ case is covered by $a=2$.

For $a=k+1$, use the crown graph $K_{a,a}$ minus a perfect matching.
It is $k$-regular, and its independence number is $a$: a mixed independent set has size at most two, while each part has size $a$.
After fewer than $a-1$ deletions, at least two vertices survive in each part and at least $a+2\ge6$ survive in total.
One part therefore has at least three vertices.
Any two vertices in the other part share a neighbor there, and every remaining vertex has a neighbor across the bipartition.
The surviving graph is connected.

The last case, $a=2$, gives an additional equality classification.

\begin{theorem}[Cycle-complement extremizers, D]\label{thm:cycles}
For $k\ge3$ and $n=k+3$, all extremizers for $f(n,2,k)$ are precisely the complements of disjoint unions of cycles of lengths at least five, with total order $n$.
\end{theorem}
\begin{proof}
Let $F$ be such a cycle union and $G=\overline F$.
Then $G$ is $(n-3)$-regular and $\alpha(G)=\omega(F)=2$.
The degree sum $n(n-3)$ is even, giving $\e(G)=nk/2$.
If fewer than $k=n-3$ deletions disconnected $G$, at least four vertices would survive.
Partition the surviving components into nonempty sets $U,V$ with no $G$-edges between them.
All cross pairs belong to $F$.
Since $F$ has maximum degree two, $|U|,|V|\le2$, so both must equal two.
Their four cross edges consume all endpoint degrees and form a whole $4$-cycle component of $F$, a contradiction.
Taking $F=C_n$ provides an example for every $n\ge6$, completing the numerical proof of Theorem~\ref{thm:shell}.

Conversely, equality and minimum degree at least $k$ force an extremizer to be $k$-regular, so its complement is $2$-regular.
It is therefore a union of cycles.
Triangles violate $\alpha(G)=2$.
For a $4$-cycle component, delete all $n-4=k-1$ vertices outside it; its complement consists of two disjoint edges, contradicting $k$-connectivity.
Every cycle must have length at least five, and the first part proves sufficiency.
\end{proof}

The constructions exhaust $2\le a\le k+1$: either $a=2$, $a\in\{k-1,k,k+1\}$, or $a\ge3$ and $k-a\ge2$.
Only the cycle-complement and tree-strip subfamilies receive complete equality classifications here.
